\frfilename{mt242.tex}
\versiondate{19.11.03}
\copyrightdate{1997}

\def\chaptername{Ruang fungsi}
\def\sectionname{$L^1$}

\newsection{242}

Meskipun ruang $L^0$ yang dibahas dalam bagian sebelumnya mempunyai daya
tarik intrinsik yang sangat besar, kegunaan utamanya dalam teori elementer
adalah sebagai ruang tempat beberapa ruang terpenting dalam analisis
fungsional dibenamkan. Dalam beberapa bagian berikutnya saya memperkenalkan
ruang-ruang tersebut satu demi satu.

Yang pertama adalah ruang $L^1$ dari kelas-kelas ekuivalensi fungsi
terintegralkan. Pentingnya ruang ini tidak hanya karena ruang ini menawarkan
bahasa untuk mengungkapkan banyak teorema mengenai fungsi terintegralkan yang
tidak bergantung pada perbedaan antara dua fungsi yang sama hampir di
mana-mana. Ruang ini juga dapat muncul sebagai ruang alami untuk mencari
penyelesaian beragam persamaan integral, dan sebagai kelengkapan suatu ruang
fungsi kontinu.

\vleader{72pt}{242A}{Ruang $L^1$} Misalkan $(X,\Sigma,\mu)$ sebarang
ruang ukur.

\header{242Aa}{\bf (a)} Misalkan $\eusm L^1=\eusm L^1(\mu)$ himpunan
fungsi-fungsi bernilai riil, yang didefinisikan pada subhimpunan-subhimpunan
$X$, dan terintegralkan pada $X$. Maka
$\eusm L^1\subseteq\eusm L^0=\eusm L^0(\mu)$\cmmnt{, sebagaimana
didefinisikan dalam \S241}, dan, untuk $f\in\eusm L^0$, kita mempunyai
$f\in\eusm L^1$ jika dan hanya jika terdapat $g\in\eusm L^1$ sedemikian sehingga
$|f|\leae g$; jika $f\in\eusm L^1$, $g\in\eusm L^0$, dan $f\eae g$,
maka $g\in\eusm L^1$. \prooflet{(Lihat 122P-122R.)}


\header{242Ab}{\bf (b)} Misalkan $L^1=L^1(\mu)\subseteq L^0=L^0(\mu)$
himpunan kelas ekuivalensi anggota-anggota $\eusm L^1$. Jika $f$,
$g\in\eusm L^1$ dan $f\eae g$, maka $\int f=\int g$\cmmnt{ (122Rb)}.
Oleh karena itu, kita dapat mendefinisikan fungsional $\int$ pada $L^1$
dengan menuliskan $\int f^{\ssbullet}=\int f$ untuk setiap
$f\in\eusm L^1$.

\header{242Ac}{\bf (c)}\cmmnt{ Akan memudahkan jika kita dapat
menuliskan} $\int_Au$ untuk $u\in L^1$, $A\subseteq X$\dvro{.}{; ini
dapat} didefinisikan dengan mengatakan bahwa
$\int_Af^{\ssbullet}=\int_Af$ untuk setiap $f\in\eusm L^1$\cmmnt{,
dengan integral tersebut didefinisikan dalam 214D}.
\prooflet{\Prf\ Saya hanya perlu memeriksa bahwa jika $f\eae g$, maka
$\int_Af=\int_Ag$; ini karena $f\restr A=g\restr A$ hampir di mana-mana
dalam $A$.\ \Qed}

Jika $E\in\Sigma$ dan $u\in L^1$, maka
$\int_Eu=\int u\times(\chi E)^{\ssbullet}$\prooflet{; ini karena
$\int_Ef=\int f\times\chi E$ untuk setiap fungsi terintegralkan $f$
(131Fa)}.

\spheader 242Ad Jika $u\in L^1$, terdapat fungsi
terukur terhadap $\Sigma$ dan terintegralkan terhadap $\mu$, yaitu $f:X\to\Bbb R$, sedemikian
sehingga $f^{\ssbullet}=u$. \prooflet{\Prf\ Sebagaimana dicatat dalam
241Bk, terdapat fungsi terukur $f:X\to\Bbb R$ sedemikian sehingga
$f^{\ssbullet}=u$; tetapi tentu saja $f$ terintegralkan karena sama hampir
di mana-mana dengan suatu fungsi terintegralkan.\ \Qed}

\leader{242B}{Teorema} Misalkan $(X,\Sigma,\mu)$ sebarang ruang ukur.
Maka $L^1(\mu)$ merupakan subruang linear dari $L^0(\mu)$ dan
$\int:L^1\to\Bbb R$ merupakan fungsional linear.

\proof{ Jika $u$, $v\in L^1=L^1(\mu)$ dan $c\in\Bbb R$, ambil
fungsi-fungsi terintegralkan $f$, $g$ sedemikian sehingga
$u=f^{\ssbullet}$ dan $v=g^{\ssbullet}$; maka $f+g$ dan $cf$
terintegralkan, sehingga $u+v=(f+g)^{\ssbullet}$ dan
$cu=(cf)^{\ssbullet}$ termasuk dalam $L^1$. Selain itu,

\Centerline{$\int u+v=\int f+g=\int f+\int g=\int u+\int v$}

\noindent dan

\Centerline{$\int cu=\int cf=c\int f=c\int u$.}
}

\leader{242C}{Struktur urutan $L^1$} Misalkan $(X,\Sigma,\mu)$
sebarang ruang ukur.

\header{242Ca}{\bf (a)} $L^1=L^1(\mu)$ mempunyai struktur urutan yang
diturunkan dari struktur urutan $L^0=L^0(\mu)$\cmmnt{ (241E); yakni,
$f^{\ssbullet}\le g^{\ssbullet}$ jika dan hanya jika $f\le g$ hampir di mana-mana}.
\cmmnt{Sebagai subruang linear dari $L^0$,} $L^1$
\dvro{merupakan}{harus merupakan} ruang linear terurut parsial\cmmnt{;
kedua syarat dalam 241Ec jelas diwarisi oleh subruang linear}.

\dvro{Jika}{Perhatikan juga bahwa jika} $u$, $v\in L^1$ dan $u\le v$,
maka $\int u\le\int v$\prooflet{, karena jika $f$, $g$ fungsi-fungsi
terintegralkan dan $f\leae g$, maka $\int f\le\int g$ (122Od)}.

\header{242Cb}{\bf (b)} Jika $u\in L^0$, $v\in L^1$, dan
$|u|\le |v|$, maka $u\in L^1$. \prooflet{\Prf\ Misalkan
$f\in\eusm L^0=\eusm L^0(\mu)$, $g\in\eusm L^1=\eusm L^1(\mu)$
sedemikian sehingga $u=f^{\ssbullet}$ dan $v=g^{\ssbullet}$; maka $g$
terintegralkan dan $|f|\leae |g|$, sehingga $f$ terintegralkan dan
$u\in L^1$.\ \Qed}

\header{242Cc}{\bf (c)} Khususnya, $|u|\in L^1$ kapan pun $u\in L^1$,
dan

\Centerline{$|\int u|=\max(\int u,\int(-u))\le\int|u|$\dvro{.}{,}}

\cmmnt{\noindent karena $u$, $-u\le|u|$.}

\header{242Cd}{\bf (d)} \cmmnt{Karena $|u|\in L^1$ untuk setiap
$u\in L^1$,

\Centerline{$u\vee v=\Bover12(u+v+|u-v|)$,
\quad$u\wedge v=\Bover12(u+v-|u-v|)$}

\noindent termasuk dalam $L^1$ untuk semua $u$, $v\in L^1$. Namun,
jika $w\in L^1$, kita tentu mempunyai

\Centerline{$w\le u\ \&\ w\le v\iff w\le u\wedge v$,}

\Centerline{$w\ge u\ \&\ w\ge v\iff w\ge u\vee v$}

\noindent sebab pernyataan-pernyataan ini berlaku untuk semua $w\in L^0$,
sehingga $u\vee v=\sup\{u,v\}$ dan $u\wedge v=\inf\{u,v\}$ dalam
$L^1$. Jadi, }$L^1$ merupakan\cmmnt{, dalam dirinya sendiri,} ruang
Riesz.

\header{242Ce}{\bf (e)} Perhatikan bahwa jika $u\in L^1$, maka
$u\ge 0$ jika dan hanya jika $\int_Eu\ge 0$ untuk setiap $E\in\Sigma$\cmmnt{; ini
karena jika $f$ suatu fungsi terintegralkan pada $X$ dan $\int_Ef\ge 0$
untuk setiap $E\in\Sigma$, maka $f\ge 0$ hampir di mana-mana (131Fb)}.
\dvro{Jika}{Secara lebih umum, jika} $u$, $v\in L^1$ dan
$\int_Eu\le\int_Ev$ untuk setiap $E\in\Sigma$, maka $u\le v$.
\dvro{Jika}{Dengan segera diperoleh bahwa jika} $u$, $v\in L^1$ dan
$\int_Eu=\int_Ev$ untuk setiap $E\in\Sigma$, maka $u=v$\cmmnt{
(bandingkan 131Fc)}.

\spheader 242Cf Jika $u\ge 0$ dalam $L^1$, terdapat fungsi nonnegatif
$f\in\eusm L^1$ sedemikian sehingga $f^{\ssbullet}=u$\cmmnt{
(bandingkan 241Eg)}.

\leader{242D}{Norma $L^1$} Misalkan $(X,\Sigma,\mu)$ sebarang ruang
ukur.

\header{242Da}{\bf (a)} Untuk $f\in\eusm L^1=\eusm L^1(\mu)$ saya
menulis $\|f\|_1=\int|f|\in\coint{0,\infty}$. Untuk
$u\in L^1=L^1(\mu)$, tetapkan $\|u\|_1=\int|u|$, sehingga
$\|f^{\ssbullet}\|_1=\|f\|_1$ untuk setiap $f\in\eusm L^1$. Maka
$\|\,\|_1$ merupakan norma pada $L^1$. \prooflet{\Prf\ (i) Jika $u$,
$v\in L^1$, maka $|u+v|\le|u|+|v|$, menurut 241Ee, sehingga

\Centerline{$\|u+v\|_1=\int|u+v|
\le\int|u|+|v|=\int|u|+\int|v|=\|u\|_1+\|v\|_1$.}

\noindent (ii) Jika $u\in L^1$ dan $c\in\Bbb R$, maka

\Centerline{$\|cu\|_1=\int|cu|=\int|c||u|=|c|\int|u|=|c|\|u\|_1$.}

\noindent (iii) Jika $u\in L^1$ dan $\|u\|_1=0$, nyatakan $u$
sebagai $f^{\ssbullet}$, dengan $f\in\eusm L^1$; maka
$\int|f|=\int|u|=0$. Karena $|f|$ nonnegatif, fungsi itu harus bernilai
nol hampir di mana-mana (122Rc), sehingga $f=0$ hampir di mana-mana dan
$u=0$ dalam $L^1$.\ \Qed}

\header{242Db}{\bf (b)}\cmmnt{ Jadi} $L^1$, dengan $\|\,\|_1$,
merupakan ruang bernorma dan $\int:L^1\to\Bbb R$ merupakan operator
linear; \cmmnt{perhatikan bahwa} $\|\int\|\le 1$\prooflet{, karena

\Centerline{$|\int u|\le\int|u|=\|u\|_1$}

\noindent untuk setiap $u\in L^1$}.

\header{242Dc}{\bf (c)} Jika $u$, $v\in L^1$ dan $|u|\le|v|$, maka

\Centerline{$\|u\|_1=\int|u|\le\int|v|=\|v\|_1$.}

\noindent Khususnya, $\|u\|_1=\||u|\|_1$ untuk setiap $u\in L^1$.

\header{242Dd}{\bf (d)} \dvro{Jika}{Perhatikan sifat berikut dari
ruang Riesz bernorma $L^1$: jika} $u$, $v\in L^1$ dan $u$, $v\ge 0$,
maka

\Centerline{$\|u+v\|_1\prooflet{\mskip5mu=\intop
u+v
=\intop u+\intop v}=\|u\|_1+\|v\|_1$.}

\spheader 242De Himpunan $(L^1)^+=\{u:u\ge 0\}$ tertutup dalam $L^1$.
\prooflet{\Prf\ Jika $v\in L^1$, $u\in(L^1)^+$, maka
$\|u-v\|_1\ge\|v\wedge 0\|_1$; ini karena jika $f$,
$g\in\eusm L^1$ dan $f\ge 0$ hampir di mana-mana,
$|f(x)-g(x)|\ge|\min(g(x),0)|$ kapan pun $f(x)$ dan $g(x)$ keduanya
terdefinisi dan $f(x)\ge 0$, yang berlaku hampir di mana-mana, sehingga

\Centerline{$\|u-v\|_1=\int|f-g|\ge\int|g\wedge\tbf{0}|
=\|v\wedge 0\|_1$.}

\noindent Kini ini berarti bahwa jika $v\in L^1$ dan $v\not\ge 0$,
bola $\{w:\|w-v\|_1<\delta\}$ tidak beririsan dengan $(L^1)^+$, dengan
$\delta=\|v\wedge 0\|_1>0$ karena $v\wedge 0\ne 0$. Jadi
$L^1\setminus(L^1)^+$ terbuka dan $(L^1)^+$ tertutup.\ \Qed}

\leader{242E}{}\cmmnt{ Untuk hasil berikutnya kita memerlukan suatu
varian teorema B.Levi.

\medskip

\noindent}{\bf Lema} Misalkan $(X,\Sigma,\mu)$ suatu ruang ukur dan
$\sequencen{f_n}$ suatu barisan fungsi riil yang terintegralkan terhadap $\mu$
sedemikian sehingga $\sum_{n=0}^{\infty}\int|f_n|<\infty$. Maka
$f=\sum_{n=0}^{\infty}f_n$ terintegralkan dan

\Centerline{$\int f=\sum_{n=0}^{\infty}\int f_n$,
\quad$\int|f|\le\sum_{n=0}^{\infty}\int|f_n|$.}

\proof{{\bf (a)} Andaikan terlebih dahulu bahwa setiap $f_n$ tak
negatif. Tetapkan $g_n=\sum_{k=0}^nf_k$ untuk setiap $n$; maka
$\sequencen{g_n}$ naik hampir di mana-mana dan

\Centerline{$\lim_{n\to\infty}\int g_n
=\sum_{k=0}^{\infty}\int f_k$}

\noindent berhingga, sehingga menurut teorema B.Levi (123A),
$f=\lim_{n\to\infty}g_n$ terintegralkan dan

\Centerline{$\int f=\lim_{n\to\infty}\int g_n
=\sum_{k=0}^{\infty}\int f_k$.}

\noindent Dalam kasus ini, tentu saja,

\Centerline{$\int|f|=\int f
=\sum_{n=0}^{\infty}\int f_n
=\sum_{n=0}^{\infty}\int|f_n|$.}

\medskip

{\bf (b)} Untuk kasus umum, tetapkan
$f_n^+={1\over 2}(|f_n|+f_n)$, $f_n^-={1\over2}(|f_n|-f_n)$, seperti
dalam 241Ef; maka $f_n^+$ dan $f_n^-$ fungsi-fungsi terintegralkan tak
negatif, dan

\Centerline{$\sum_{n=0}^{\infty}\int f_n^+
  +\sum_{n=0}^{\infty}\int f_n^-
=\sum_{n=0}^{\infty}\int|f_n|<\infty$.}

\noindent Jadi $h_1=\sum_{n=0}^{\infty}f_n^+$ dan
$h_2=\sum_{n=0}^{\infty}f_n^-$ keduanya terintegralkan. Kini
$f\eae h_1-h_2$, sehingga

\Centerline{$\int f=\int h_1-\int h_2
=\sum_{n=0}^{\infty}\int f_n^+
  -\sum_{n=0}^{\infty}\int f_n^-
=\sum_{n=0}^{\infty}\int f_n$.}

\noindent Akhirnya,

\Centerline{$\int|f|\le\int|h_1|+\int|h_2|
=\sum_{n=0}^{\infty}\int f_n^+
  +\sum_{n=0}^{\infty}\int f_n^-
=\sum_{n=0}^{\infty}\int|f_n|$.}
}

\leader{242F}{Teorema} Untuk sebarang ruang ukur $(X,\Sigma,\mu)$,
$L^1(\mu)$ lengkap terhadap normanya $\|\,\|_1$.

\proof{ Misalkan $\sequencen{u_n}$ suatu barisan dalam $L^1$ sedemikian
sehingga $\|u_{n+1}-u_n\|_1\le 4^{-n}$ untuk setiap $n\in\Bbb N$.
Pilih fungsi-fungsi terintegralkan $f_n$ sedemikian sehingga
$f_0^{\ssbullet}=u_0$, $f_{n+1}^{\ssbullet}=u_{n+1}-u_n$ untuk setiap
$n\in\Bbb N$. Maka

\Centerline{$\sum_{n=0}^{\infty}\int|f_n|
=\|u_0\|_1+\sum_{n=0}^{\infty}\|u_{n+1}-u_n\|_1<\infty$.}

\noindent Jadi $f=\sum_{n=0}^{\infty}f_n$ terintegralkan menurut 242E,
dan $u=f^{\ssbullet}\in L^1$. Tetapkan $g_n=\sum_{j=0}^nf_j$ untuk
setiap $n$; maka $g_n^{\ssbullet}=u_n$, sehingga

\Centerline{$\|u-u_n\|_1
=\int|f-g_n|
\le\int\sum_{j=n+1}^{\infty}|f_j|
\le\sum_{j=n+1}^{\infty}4^{-j}
=4^{-n}/3$}

\noindent untuk setiap $n$. Jadi $u=\lim_{n\to\infty}u_n$ dalam
$L^1$. Karena $\sequencen{u_n}$ sebarang, $L^1$ lengkap (2A4E).
}

\leader{242G}{Definisi}\cmmnt{ Untuk rujukan kelak, akan memudahkan
jika kita memperkenalkan istilah berikut.} Suatu {\bf kisi Banach}
adalah ruang Riesz $U$ bersama suatu norma $\|\,\|$ pada $U$ sedemikian
sehingga (i) $\|u\|\le\|v\|$ kapan pun $u$, $v\in U$ dan
$|u|\le|v|$\cmmnt{, dengan $|u|$ menyatakan $u\vee(-u)$, seperti dalam
241Ee} (ii) $U$ lengkap terhadap $\|\,\|$. \cmmnt{Jadi 242Dc dan 242F
secara bersama-sama menyatakan bahwa ruang Riesz bernorma
$(L^1,\|\,\|_1)$ merupakan kisi Banach.}

\leader{242H}{$L^1$ sebagai \dvrocolon{ruang Riesz}}\cmmnt{ Kita dapat
membahas ruang linear terurut $L^1$ dengan bahasa yang telah digunakan
dalam 241E-241G %241E 241F 241G
untuk $L^0$.

\medskip

\noindent}{\bf Teorema} Misalkan $(X,\Sigma,\mu)$ sebarang ruang ukur.
Maka $L^1=L^1(\mu)$ lengkap Dedekind.
\proof{{\bf (a)} Misalkan $A\subseteq L^1$ sebarang himpunan tak kosong
yang terbatas di atas dalam $L^1$. Tetapkan

\Centerline{$A'=\{u_0\vee\ldots\vee u_n:u_0,\ldots, u_n\in A\}$.}

\noindent Maka $A\subseteq A'$, $A'$ mempunyai batas-batas atas yang
sama dengan $A$, dan $u\vee v\in A'$ untuk semua $u$, $v\in A'$.
Dengan mengambil $w_0$ sebagai sebarang batas atas $A$ dan $A'$, kita
mempunyai $\int u\le\int w_0$ untuk setiap $u\in A'$, sehingga
$\gamma=\sup_{u\in A'}\int u$ terdefinisi dalam $\Bbb R$. Untuk setiap
$n\in\Bbb N$, pilih $u_n\in A'$ sedemikian sehingga
$\int u_n\ge\gamma-2^{-n}$. Karena $L^0=L^0(\mu)$ lengkap-$\sigma$
Dedekind (241Ga), $u^*=\sup_{n\in\Bbb N}u_n$ terdefinisi dalam $L^0$,
dan $u_0\le u^*\le w_0$ dalam $L^0$. Akibatnya,

\Centerline{$0\le u^*-u_0\le w_0-u_0$}

\noindent dalam $L^0$. Namun $w_0-u_0\in L^1$, sehingga
$u^*-u_0\in L^1$ (242Cb) dan $u^*\in L^1$.

\medskip

{\bf (b)} Pokoknya adalah bahwa $u^*$ merupakan batas atas bagi $A$.
\Prf\ Jika $u\in A$, maka $u\vee u_n\in A'$ untuk setiap $n$, sehingga

$$\eqalignno{\|u-u\wedge u^*\|_1
&=\int u-u\wedge u^*
\le \int u-u\wedge u_n\cr
\noalign{\noindent (karena $u\wedge u_n\le u_n\le u^*$, sehingga
$u\wedge u_n\le u\wedge u^*$)}
&=\int u\vee u_n-u_n\cr
\noalign{\noindent (karena $u\vee u_n+u\wedge u_n=u+u_n$ -- lihat
rumus-rumus dalam 242Cd)}
&=\int u\vee u_n-\int u_n
\le\gamma-(\gamma-2^{-n})
=2^{-n}\cr}$$

\noindent untuk setiap $n$; jadi $\|u-u\wedge u^*\|_1=0$. Namun ini
berarti bahwa $u=u\wedge u^*$, yakni $u\le u^*$. Karena $u$ sebarang,
$u^*$ merupakan batas atas bagi $A$.\ \Qed

\medskip

{\bf (c)} Di sisi lain, setiap batas atas bagi $A$ tentu merupakan batas
atas bagi $\{u_n:n\in\Bbb N\}$, sehingga lebih besar daripada atau sama
dengan $u^*$. Jadi $u^*=\sup A$ dalam $L^1$. Karena $A$ sebarang,
$L^1$ lengkap Dedekind.
}%end of proof of 242H

\cmmnt{
\medskip

\noindent{\bf Catatan} Perhatikan bahwa kelengkapan menurut urutan dari
$L^1$, tidak seperti kelengkapan menurut urutan dari $L^0$, tidak bergantung
pada sifat khusus apa pun dari ruang ukur $(X,\Sigma,\mu)$.
}%end of comment

\leader{242I}{Teorema Radon-Nikod\'ym \dvro{}{}}\cmmnt{ Menurut saya,
teorema Radon-Nikod\'ym (232E) layak dituliskan kembali dalam bahasa bab ini.

\medskip

\noindent{\bf Teorema}} Misalkan $(X,\Sigma,\mu)$ suatu ruang ukur.
Maka terdapat bijeksi kanonik antara $L^1=L^1(\mu)$ dan himpunan fungsional
aditif kontinu sejati $\nu:\Sigma\to\Bbb R$, yang diberikan oleh rumus

\Centerline{$\nu F=\int_Fu$ untuk $F\in\Sigma$, $u\in L^1$.}

\cmmnt{\medskip

\noindent{\bf Catatan} Ingat bahwa jika $\mu$ $\sigma$-hingga, maka
fungsional aditif kontinu sejati tepat merupakan fungsional aditif terhitung
yang kontinu mutlak; dan jika $\mu$ hingga total, maka semua fungsional
aditif (hingga) yang kontinu mutlak bersifat kontinu sejati (232Bd).
}%end of comment

\proof{ Untuk $u\in L^1$, $F\in\Sigma$, tetapkan
$\nu_uF=\int_Fu$. Jika $u\in L^1$, terdapat fungsi terintegralkan $f$
sedemikian sehingga $f^{\ssbullet}=u$, dan dalam hal ini

\Centerline{$F\mapsto\nu_uF=\int_Ff:\Sigma\to\Bbb R$}

\noindent bersifat aditif dan kontinu sejati menurut 232D. Jika
$\nu:\Sigma\to\Bbb R$ aditif dan kontinu sejati, maka menurut 232E
terdapat fungsi terintegralkan $f$ sedemikian sehingga
$\nu F=\int_Ff$ untuk setiap $F\in\Sigma$; dengan menetapkan
$u=f^{\ssbullet}$ dalam $L^1$, diperoleh $\nu=\nu_u$. Akhirnya, jika
$u$, $v$ merupakan anggota berbeda dari $L^1$, terdapat $F\in\Sigma$
sedemikian sehingga $\int_Fu\ne\int_Fv$ (242Ce), sehingga
$\nu_u\ne\nu_v$; jadi $u\mapsto\nu_u$ injektif sekaligus surjektif.
}%end of proof of 242I

\leader{242J}{Meninjau kembali ekspektasi bersyarat}\cmmnt{ Kini kita
mempunyai perangkat yang diperlukan untuk memberikan penafsiran baru atas
beberapa gagasan dalam \S233.

\medskip

}{\bf (a)} Misalkan $(X,\Sigma,\mu)$ suatu ruang ukur, dan $\Tau$ suatu
subaljabar-$\sigma$ dari $\Sigma$\cmmnt{, seperti dalam 233A}. Maka
$(X,\Tau,\mu\restrp\Tau)$ suatu ruang ukur, dan
$\eusm L^0(\mu\restrp\Tau)\subseteq\eusm L^0(\mu)$; \cmmnt{selain itu,}
jika $f$, $g\in\eusm L^0(\mu\restrp\Tau)$, maka
$f=g\,\,(\mu\restrp\Tau)$-hampir di mana-mana jika dan hanya jika
$f=g\,\,\mu$-hampir di mana-mana. \prooflet{\Prf\ Terdapat himpunan-himpunan
yang koterabaikan terhadap $\mu\restrp\Tau$, yaitu $F$, $G\in\Tau$, sedemikian
sehingga $f\restr F$ dan $g\restr G$ terukur terhadap $\Tau$; tetapkan

\Centerline{$E=\{x:x\in F\cap G,\,f(x)\ne g(x)\}\in\Tau$;}

\noindent maka

\Centerline{$f=g\,\,(\mu\restrp\Tau)$-hampir di mana-mana $\iff
(\mu\restrp\Tau)(E)=0
\iff \mu E=0 \iff f=g\,\,\mu$-hampir di mana-mana. \Qed}
}

Oleh karena itu, kita mempunyai peta kanonik
$S:L^0(\mu\restrp\Tau)\to L^0(\mu)$ yang didefinisikan dengan mengatakan
bahwa jika $u\in L^0(\mu\restrp\Tau)$ merupakan kelas ekuivalensi dari
$f\in\eusm L^0(\mu\restrp\Tau)$, maka $Su$ merupakan kelas ekuivalensi
dari $f$ dalam $L^0(\mu)$. \cmmnt{Dengan menelusuri operasi-operasi yang
dideskripsikan dalam 241D, 241E, dan 241H, mudah diperiksa bahwa} $S$ linear,
injektif, dan mempertahankan urutan, serta\cmmnt{ bahwa}
$|Su|=S|u|$, $S(u\vee v)=Su\vee Sv$, dan
$S(u\times v)=Su\times Sv$ untuk $u$, $v\in L^0(\mu\restrp\Tau)$.

\header{242Jb}{\bf (b)} Selanjutnya, jika
$f\in\eusm L^1(\mu\restrp\Tau)$, maka $f\in\eusm L^1(\mu)$ dan
$\int fd\mu=\int fd(\mu\restrp\Tau)$\prooflet{ (233B)}; jadi
$Su\in L^1(\mu)$ dan $\|Su\|_1=\|u\|_1$ untuk setiap
$u\in L^1(\mu\restrp\Tau)$.

Perhatikan juga bahwa setiap anggota dari
$L^1(\mu)\cap S[L^0(\mu\restrp\Tau)]$ sesungguhnya termasuk dalam
$S[L^1(\mu\restrp\Tau)]$. \prooflet{\Prf\ Ambil
$u\in L^1(\mu)\cap S[L^0(\mu\restrp\Tau)]$. Maka $u$ dapat dinyatakan
baik sebagai $f^{\ssbullet}$ dengan $f\in\eusm L^1(\mu)$, maupun sebagai
$g^{\ssbullet}$ dengan $g\in\eusm L^0(\mu\restrp\Tau)$. Jadi
$g\eae f$, dan $g$ terintegralkan terhadap $\mu$, sehingga terintegralkan
terhadap $(\mu\restrp\Tau)$ (sekali lagi menurut 233B).\ \Qed}

Ini berarti bahwa
$S:L^1(\mu\restrp\Tau)\to L^1(\mu)\cap S[L^0(\mu\restrp\Tau)]$
merupakan bijeksi.

\header{242Jc}{\bf (c)} Kini andaikan $\mu X=1$, sehingga
$(X,\Sigma,\mu)$ suatu ruang probabilitas. \cmmnt{Ingat bahwa $g$
merupakan ekspektasi bersyarat dari $f$ pada $\Tau$ jika $g$
terintegralkan terhadap $\mu\restrp\Tau$, $f$ terintegralkan terhadap
$\mu$, dan $\int_Fg=\int_Ff$ untuk setiap $F\in\Tau$; serta bahwa setiap
fungsi yang terintegralkan terhadap $\mu$ mempunyai ekspektasi bersyarat semacam itu
(233D).} Jika $g$ merupakan ekspektasi bersyarat dari $f$ dan
$f_1=f\,\,\mu$-hampir di mana-mana, maka $g$ merupakan ekspektasi
bersyarat dari $f_1$\cmmnt{, karena $\int_Ff_1=\int_Ff$ untuk setiap
$F$}; dan \cmmnt{saya telah mencatat dalam 233Dc bahwa} jika $g$,
$g_1$ merupakan ekspektasi bersyarat dari $f$ pada $\Tau$, maka
$g=g_1\,\,\mu\restrp\Tau$-hampir di mana-mana.

\header{242Jd}{\bf (d)} Ini berarti bahwa kita mempunyai operator
$P:L^1(\mu)\to L^1(\mu\restrp\Tau)$ yang didefinisikan dengan mengatakan
bahwa $P(f^{\ssbullet})=g^{\ssbullet}$ kapan pun
$g\in\eusm L^1(\mu\restrp\Tau)$ merupakan ekspektasi bersyarat dari
$f\in\eusm L^1(\mu)$ pada $\Tau$\cmmnt{; yakni,
$\int_FPu=\int_Fu$ kapan pun $u\in L^1(\mu)$ dan $F\in\Tau$}.
\cmmnt{Jika kita mengidentifikasi $L^1(\mu)$,
$L^1(\mu\restrp\Tau)$ dengan himpunan fungsional aditif kontinu mutlak
yang didefinisikan pada $\Sigma$ dan $\Tau$, seperti dalam 242I, maka
$P$ bersesuaian dengan operasi $\nu\mapsto\nu\restrp\Tau$.}

\header{242Je}{\bf (e)} \dvro{$P$ linear dan mempertahankan urutan.}
{Karena $Pu$ didefinisikan secara unik dalam $L^1(\mu\restrp\Tau)$ oleh
syarat $\int_FPu=\int_Fu$ untuk setiap $F\in\Tau$ (242Ce), kita melihat
bahwa $P$ harus linear.} \prooflet{\Prf\ Jika $u$, $v\in L^1(\mu)$ dan
$c\in\Bbb R$, maka

\Centerline{$\int_FPu+Pv=\int_FPu+\int_FPv=\int_Fu+\int_Fv
=\int_Fu+v=\int_FP(u+v)$,}

\Centerline{$\int_FP(cu)=\int_Fcu=c\int_Fu=c\int_FPu=\int_FcPu$}

\noindent untuk setiap $F\in\Tau$.\ \QeD\ Selain itu, jika $u\ge 0$,
maka $\int_FPu=\int_Fu\ge 0$ untuk setiap $F\in\Tau$, sehingga
$Pu\ge 0$ (sekali lagi menurut 242Ce).

Dengan segera diperoleh bahwa $P$ mempertahankan urutan, yakni
$Pu\le Pv$ kapan pun $u\le v$.} Akibatnya,

\Centerline{$|Pu|=Pu\vee(-Pu)=Pu\vee P(-u)\le P|u|$}

\noindent untuk setiap $u\in L^1(\mu)$\prooflet{, karena
$u\le|u|$ dan $-u\le|u|$}. Akhirnya, $P$ merupakan operator linear
terbatas, dengan norma $1$. \prooflet{\Prf\ Rumus terakhir menyatakan
bahwa

\Centerline{$\|Pu\|_1\le\|P|u|\|_1=\int P|u|=\int |u|=\|u\|_1$}

\noindent untuk setiap $u\in L^1(\mu)$, sehingga $\|P\|\le 1$. Di sisi
lain, $P(\chi X^{\ssbullet})=\chi X^{\ssbullet}\ne 0$, sehingga
$\|P\|=1$.\ \Qed}

\header{242Jf}{\bf (f)} Kita dapat secara sah memandang
$Pu\in L^1(\mu\restrp\Tau)$ sebagai ekspektasi bersyarat `dari'
$u\in L^1(\mu)$ pada $\Tau$; $P$ merupakan {\bf operator ekspektasi
bersyarat}.

\header{242Jg}{\bf (g)} Jika $u\in L^1(\mu\restrp\Tau)$, maka\cmmnt{
kita mempunyai} $Su\in L^1(\mu)$ yang bersesuaian, seperti dalam (b);
kini $PSu=u$. \prooflet{\Prf\
$\int_FPSu=\int_FSu=\int_Fu$ untuk setiap $F\in\Tau$.\ \QeD}
Akibatnya, $SPSP=SP:L^1(\mu)\to L^1(\mu)$.

\cmmnt{\header{242Jh}{\bf (h)} Pembedaan yang dibuat di atas antara
$u=f^{\ssbullet}\in L^0(\mu\restrp\Tau)$ dan
$Su=f^{\ssbullet}\in L^0(\mu)$ tentu saja sangat teliti. Saya percaya
kita perlu menyadari pembedaan semacam itu, meskipun untuk hampir semua
keperluan aman sekaligus mudah untuk memandang $L^0(\mu\restrp\Tau)$
sebagai subhimpunan aktual dari $L^0(\mu)$. Jika kita melakukannya, maka
(b) menyatakan bahwa kita dapat mengidentifikasi $L^1(\mu\restrp\Tau)$
dengan $L^1(\mu)\cap L^0(\mu\restrp\Tau)$, sedangkan (g) menjadi
`$P^2=P$'.
}

\leader{242K}{}\cmmnt{ Bahasa yang baru diperkenalkan memungkinkan
perumusan ulang 233J-233K berikut.

\medskip

\noindent}{\bf Teorema} Misalkan $(X,\Sigma,\mu)$ suatu ruang
probabilitas dan $\Tau$ suatu subaljabar-$\sigma$ dari $\Sigma$.
Misalkan $\phi:\Bbb R\to\Bbb R$ suatu fungsi cembung dan
$\bar\phi:L^0(\mu)\to L^0(\mu)$ operator yang bersesuaian, yang
didefinisikan dengan menetapkan
$\bar\phi(f^{\ssbullet})=(\phi f)^{\ssbullet}$\cmmnt{ (241I)}. Jika
$P:L^1(\mu)\to L^1(\mu\restrp\Tau)$ merupakan operator ekspektasi
bersyarat, maka $\bar\phi(Pu)\le P(\bar\phi u)$ kapan pun
$u\in L^1(\mu)$ sedemikian sehingga $\bar\phi(u)\in L^1(\mu)$.

\proof{ Ini hanyalah pernyataan ulang 233J.}

\leader{242L}{Proposisi} Misalkan $(X,\Sigma,\mu)$ suatu ruang
probabilitas, dan $\Tau$ suatu subaljabar-$\sigma$ dari $\Sigma$.
Misalkan $P:L^1(\mu)\to L^1(\mu\restrp\Tau)$ operator ekspektasi
bersyarat yang bersesuaian. Jika $u\in L^1=L^1(\mu)$ dan
$v\in L^0(\mu\restrp\Tau)$, maka $u\times v\in L^1$ jika dan hanya jika
$P|u|\times v\in L^1$, dan dalam hal ini
$P(u\times v)=Pu\times v$; khususnya,
$\int u\times v=\int Pu\times v$.

\proof{ (Di sini saya menggunakan identifikasi $L^0(\mu\restrp\Tau)$
sebagai subruang dari $L^0(\mu)$, sebagaimana disarankan dalam 242Jh.)
Nyatakan $u$ sebagai $f^{\ssbullet}$ dan $v$ sebagai $h^{\ssbullet}$,
dengan $f\in\eusm L^1=\eusm L^1(\mu)$ dan
$h\in\eusm L^0(\mu\restrp\Tau)$. Misalkan
$g$, $g_0\in\eusm L^1(\mu\restrp\Tau)$ berturut-turut merupakan
ekspektasi bersyarat dari $f$, $|f|$, sehingga $Pu=g^{\ssbullet}$ dan
$P|u|=g_0^{\ssbullet}$. Kemudian, dengan menggunakan 233K,

\Centerline{$u\times v\in L^1\iff f\times h\in\eusm L^1
\iff g_0\times h\in\eusm L^1\iff P|u|\times v\in L^1$,}

\noindent dan dalam hal ini $g\times h$ merupakan ekspektasi bersyarat
dari $f\times h$, yakni $Pu\times v=P(u\times v)$.
}%end of proof of 242L

\leader{242M}{$L^1$ sebagai \dvro{kelengkapan:}{kelengkapan}}\cmmnt{
Dalam pendahuluan bagian ini saya menyebutkan bahwa $L^1$ muncul dalam
analisis fungsional sebagai kelengkapan beberapa ruang penting; dengan kata
lain, beberapa subruang padat dari $L^1$ memiliki arti penting. Yang pertama
bersifat elementer.

\medskip

\noindent}{\bf Proposisi} Misalkan $(X,\Sigma,\mu)$ sebarang ruang ukur,
dan tuliskan $\eusm S$ untuk ruang fungsi-fungsi sederhana terhadap $\mu$
pada $X$. Maka

(a) kapan pun $f$ suatu fungsi bernilai riil yang terintegralkan terhadap $\mu$ dan
$\epsilon>0$, terdapat $h\in\eusm S$ sedemikian sehingga
$\int|f-h|\le\epsilon$;

(b) $S=\{f^{\ssbullet}:f\in\eusm S\}$ merupakan subruang linear padat
dari $L^1=L^1(\mu)$.

\proof{{\bf (a)}(i) Jika $f$ nonnegatif, terdapat fungsi sederhana $h$
sedemikian sehingga $h\leae f$ dan
$\int h\ge\int f-\bover12\epsilon$ (122K), dan dalam hal ini

\Centerline{$\int|f-h|=\int f-h
=\int f-\int h\le\Bover12\epsilon$.}

\noindent (ii) Dalam kasus umum, $f$ dapat dinyatakan sebagai selisih
$f_1-f_2$ dari fungsi-fungsi terintegralkan nonnegatif. Kini terdapat
$h_1$, $h_2\in\eusm S$ sedemikian sehingga
$\int|f_j-h_j|\le\bover12\epsilon$ untuk kedua nilai $j$, dan

\Centerline{$\int|f-h|\le\int|f_1-h_1|+\int|f_2-h_2|\le\epsilon$.}

\medskip

{\bf (b)} Karena $\eusm S$ merupakan subruang linear dari $\Bbb R^X$
yang termuat dalam $\eusm L^1=\eusm L^1(\mu)$, $S$ merupakan subruang
linear dari $L^1$. Jika $u\in L^1$ dan $\epsilon>0$, terdapat
$f\in\eusm L^1$ sedemikian sehingga $f^{\ssbullet}=u$ dan
$h\in\eusm S$ sedemikian sehingga $\int|f-h|\le\epsilon$; kini
$v=h^{\ssbullet}\in S$ dan

\Centerline{$\|u-v\|_1=\int|f-h|\le\epsilon$.}

\noindent Karena $u$ dan $\epsilon$ sebarang, $S$ padat dalam $L^1$.
}

\leader{242N}{}\cmmnt{ Seperti biasa, ukuran Lebesgue pada $\BbbR^r$
dan subhimpunan-subhimpunannya sejauh ini merupakan contoh terpenting; dalam
kasus ini kita mempunyai kelas-kelas lain dari subruang padat $L^1$. Jika
Anda telah sampai di titik ini tanpa berusaha menguasai ukuran Lebesgue
multidimensi, ambil saja $r=1$. Jika Anda kurang nyaman dengan ukuran
subruang umum, ambil $X$ sebagai $\BbbR^r$ atau $[0,1]\subseteq\Bbb R$
atau subhimpunan tertentu lain yang Anda anggap menarik. Istilah berikut
akan berguna.

\medskip

\noindent}{\bf Definisi} Jika $f$ suatu fungsi bernilai riil atau kompleks
yang didefinisikan pada subhimpunan $\BbbR^r$, katakan bahwa
{\bf dukungan} $f$ adalah
$\overline{\{x:x\in\dom f,\,f(x)\ne 0\}}$.

\leader{242O}{Teorema} Misalkan $X$ sebarang subhimpunan $\BbbR^r$,
dengan $r\ge 1$, dan misalkan $\mu$ ukuran Lebesgue pada $X$\cmmnt{,
yakni ukuran subruang pada $X$ yang diimbaskan oleh ukuran Lebesgue pada
$\BbbR^r$}. Tuliskan $C_k$ untuk ruang fungsi kontinu terbatas
$f:\BbbR^r\to\Bbb R$ yang mempunyai dukungan terbatas, dan $S_0$ untuk
ruang kombinasi linear fungsi-fungsi berbentuk $\chi I$, dengan
$I\subseteq\BbbR^r$ suatu interval setengah terbuka terbatas. Maka

(a) kapan pun $f\in\eusm L^1=\eusm L^1(\mu)$ dan $\epsilon>0$,
terdapat $g\in C_k$, $h\in S_0$ sedemikian sehingga
$\int_X|f-g|\le\epsilon$ dan $\int_X|f-h|\le\epsilon$;

(b) $\{(g\restr X)^{\ssbullet}:g\in C_k\}$ dan
$\{(h\restr X)^{\ssbullet}:h\in S_0\}$ merupakan subruang-subruang
linear padat dari $L^1=L^1(\mu)$.

\cmmnt{\medskip

\noindent{\bf Catatan} Tentu saja kata `terbatas' dalam deskripsi $C_k$
bersifat berlebih; lihat 242Xh.}

\proof{{\bf (a)} Saya menunjukkan secara berturut-turut bahwa hasil ini
berlaku bagi makin banyak anggota $f$ dari
$\eusm L^1=\eusm L^1(\mu)$. Tuliskan $\mu_r$ untuk ukuran Lebesgue pada
$\BbbR^r$, sehingga $\mu$ merupakan ukuran subruang $(\mu_r)_X$.

\medskip

\quad{\bf (i)} Andaikan terlebih dahulu bahwa $f=\chi I\restr X$ dengan
$I\subseteq\BbbR^r$ suatu interval setengah terbuka terbatas. Tentu saja
$\chi I$ sudah termasuk dalam $S_0$, sehingga saya hanya perlu menunjukkan
bahwa fungsi itu dapat dihampiri oleh anggota-anggota $C_k$. Jika
$I=\emptyset$, hasilnya sepele; kita dapat mengambil $g=0$. Jika tidak,
nyatakan $I$ sebagai $\coint{a-b,a+b}$, dengan
$a=(\alpha_1,\ldots,\alpha_r)$, $b=(\beta_1,\ldots,\beta_r)$, dan
$\beta_j>0$ untuk setiap $j$. Ambil $\delta>0$ sedemikian sehingga

\Centerline{$2^r\prod_{j=1}^r(\beta_j+\delta)
\le\epsilon+2^r\prod_{j=1}^r\beta_j$.}

\noindent Untuk $\xi\in\Bbb R$, tetapkan

$$\eqalign{g_j(\xi)
&=1\text{ jika }|\xi-\alpha_j|\le\beta_j,\cr
&=(\beta_j+\delta-|\xi-\alpha_j|)/\delta
  \text{ jika }\beta_j\le|\xi-\alpha_j|\le\beta_j+\delta,\cr
&=0\text{ jika }|\xi-\alpha_j|\ge\beta_j+\delta.\cr}$$

\def\Caption{Fungsi $g_j$}
\picture{mt242m}{100pt}

\noindent Untuk $x=(\xi_1,\ldots,\xi_r)\in\BbbR^r$, tetapkan

\Centerline{$g(x)=\prod_{j=1}^rg_j(\xi_j)$.}

\noindent Maka $g\in C_k$ dan $\chi I\le g\le\chi J$, dengan
$J=[a-b-\delta \tbf{1},a+b+\delta\tbf{1}]$ (dengan menuliskan
$\tbf{1}=(1,\ldots,1)$), sehingga (berdasarkan pilihan $\delta$)
$\mu_rJ\le\mu_rI+\epsilon$, dan

$$\eqalign{\int_X|g-f|
&\le\int(\chi(J\cap X)-\chi(I\cap X))d\mu
=\mu((J\setminus I)\cap X)\cr
&\le\mu_r(J\setminus I)
=\mu_rJ-\mu_rI
\le\epsilon,\cr}$$

\noindent sebagaimana diperlukan.

\medskip

\quad{\bf (ii)} Kini andaikan $f=\chi(X\cap E)$ dengan
$E\subseteq\BbbR^r$ suatu himpunan berukuran hingga. Maka terdapat keluarga
terpisah $I_0,\ldots,I_n$ dari interval-interval setengah terbuka sedemikian
sehingga $\mu_r(E\symmdiff\bigcup_{j\le n}I_j)\le\bover12\epsilon$.
\Prf\ Terdapat himpunan terbuka $G\supseteq E$ sedemikian sehingga
$\mu_r(G\setminus E)\le\bover14\epsilon$ (134Fa). Untuk setiap
$m\in\Bbb N$, misalkan $\Cal I_m$ keluarga interval setengah terbuka dalam
$\BbbR^r$ berbentuk $\coint{a,b}$, dengan
$a=(2^{-m}k_1,\ldots,2^{-m}k_r)$, $k_1,\ldots,k_r$ bilangan bulat, dan
$b=a+2^{-m}\tbf{1}$; maka $\Cal I_m$ merupakan keluarga terpisah.
Tetapkan $H_m=\bigcup\{I:I\in\Cal I_m,\,I\subseteq G\}$; maka
$\sequence{m}{H_m}$ merupakan keluarga tak menurun dengan gabungan $G$,
sehingga terdapat $m$ sedemikian sehingga
$\mu_r(G\setminus H_m)\le\bover14\epsilon$ dan
$\mu_r(E\symmdiff H_m)\le\bover12\epsilon$. Namun kini $H_m$ dapat
dinyatakan sebagai gabungan terpisah $\bigcup_{j\le n}I_j$, dengan
$I_0,\ldots,I_n$ mencacah anggota-anggota $\Cal I_m$ yang termuat dalam
$H_m$. (Kalimat terakhir gagal jika $H_m$ kosong. Namun jika
$H_m=\emptyset$, kita dapat mengambil $n=0$, $I_0=\emptyset$.)\ \Qed

Oleh karena itu, $h=\sum_{j=0}^n\chi I_j\in S_0$ dan

\Centerline{$\int_X|f-h|
=\mu(X\cap(E\symmdiff\bigcup_{j\le n}I_j))\le\Bover12\epsilon$.}

\noindent Untuk $C_k$, butir (i) menyatakan bahwa untuk setiap $j\le n$
terdapat $g_j\in C_k$ sedemikian sehingga
$\int_X|g_j-\chi I_j|\le\epsilon/2(n+1)$, sehingga
$g=\sum_{j=0}^ng_j\in C_k$ dan

\Centerline{$\int_X|f-g|\le\int_X|f-h|+\int_X|h-g|
\le\Bover{\epsilon}2+\sum_{j=0}^n\int_X|g_j-\chi I_j|\le\epsilon$.}

\medskip

\quad{\bf (iii)} Jika $f$ suatu fungsi sederhana, nyatakan $f$ sebagai
$\sum_{k=0}^na_k\chi E_k$, dengan setiap $E_k$ berukuran hingga terhadap
$\mu$. Setiap $E_k$ dapat dinyatakan sebagai $X\cap F_k$, dengan
$\mu_rF_k=\mu E_k$ (214Ca). Menurut (ii), kita dapat menemukan
$g_k\in C_k$, $h_k\in S_0$ sedemikian sehingga

\Centerline{$|a_k|\int_X|g_k-\chi F_k|\le\Bover{\epsilon}{n+1}$,
\quad$|a_k|\int_X|h_k-\chi F_k|\le\Bover{\epsilon}{n+1}$}

\noindent untuk setiap $k$. Tetapkan $g=\sum_{k=0}^na_kg_k$ dan
$h=\sum_{k=0}^na_kh_k$; maka $g\in C_k$, $h\in S_0$, dan

\Centerline{$\int_X|f-g|\le\int_X\sum_{k=0}^n|a_k||\chi F_k-g_k|
=\sum_{k=0}^n|a_k|\int_X|\chi F_k-g_k|\le\epsilon$,}

\Centerline{$\int_X|f-h|
\le\sum_{k=0}^n|a_k|\int_X|\chi F_k-h_k|\le\epsilon$,}

\noindent sebagaimana diperlukan.

\medskip

\quad{\bf (iv)} Jika $f$ sebarang fungsi terintegralkan pada $X$, maka
menurut 242Ma kita dapat menemukan fungsi sederhana $f_0$ sedemikian
sehingga $\int|f-f_0|\le\bover12\epsilon$, dan menurut (iii) terdapat
$g\in C_k$, $h\in S_0$ sedemikian sehingga
$\int_X|f_0-g|\le\bover12\epsilon$,
$\int_X|f_0-h|\le\bover12\epsilon$; sehingga

\Centerline{$\int_X|f-g|
\le\int_X|f-f_0|+\int_X|f_0-g|\le\epsilon$,}

\Centerline{$\int_X|f-h|
\le\int_X|f-f_0|+\int_X|f_0-h|\le\epsilon$.}


\medskip

{\bf (b)(i)} Pertama-tama kita harus memeriksa bahwa jika $g\in C_k$,
maka $g\restr X$ memang terintegralkan terhadap $\mu$. Pokoknya di sini adalah
bahwa jika $g\in C_k$ dan $a\in\Bbb R$, maka

\Centerline{$\{x:x\in X,\,g(x)>a\}$}

\noindent merupakan irisan $X$ dengan suatu subhimpunan terbuka dari
$\BbbR^r$, dan karenanya terukur oleh $\mu$, sebab semua himpunan terbuka
terukur oleh $\mu_r$ (115G). Selanjutnya, $g$ terbatas dan himpunan
$E=\{x:x\in X,\,g(x)\ne 0\}$ terbatas dalam $\BbbR^r$, sehingga
berukuran luar hingga terhadap $\mu_r$ dan berukuran hingga terhadap
$\mu$. Jadi, terdapat $M\ge 0$ sedemikian sehingga
$|g|\le M\chi E$, yang terintegralkan terhadap $\mu$. Oleh karena itu, $g$
terintegralkan terhadap $\mu$.

Tentu saja $h\restr X$ terintegralkan terhadap $\mu$ untuk setiap $h\in S_0$,
karena (berdasarkan definisi ukuran subruang) $\mu(I\cap X)$ terdefinisi
dan berhingga untuk setiap interval setengah terbuka terbatas $I$.
\medskip

\quad{\bf (ii)} Kini sisanya mengikuti dengan argumen yang persis sama
seperti dalam 242Mb. Karena $\{g\restr X:g\in C_k\}$ dan
$\{h\restr X:h\in S_0\}$ merupakan subruang-subruang linear dari
$\Bbb R^X$ yang termuat dalam $\eusm L^1(\mu)$, citra-citranya
$C_k^{\#}$ dan $S_0^{\#}$ merupakan subruang-subruang linear dari $L^1$.
Jika $u\in L^1$ dan $\epsilon>0$, terdapat $f\in\eusm L^1$ sedemikian
sehingga $f^{\ssbullet}=u$, serta $g\in C_k$, $h\in S_0$ sedemikian
sehingga $\int_X|f-g|$, $\int_X|f-h|\le\epsilon$; kini
$v=(g\restr X)^{\ssbullet}\in C_k^{\#}$ dan
$w=(h\restr X)^{\ssbullet}\in S_0^{\#}$, dan

\Centerline{$\|u-v\|_1=\int_X|f-g|\le\epsilon$,
\quad$\|u-w\|_1=\int_X|f-h|\le\epsilon$.}

\noindent Karena $u$ dan $\epsilon$ sebarang, $C_k^{\#}$ dan
$S_0^{\#}$ padat dalam $L^1$.
}%end of proof of 242O

\leader{242P}{$L^1$ kompleks}\cmmnt{ Sebagaimana, saya harap, Anda
duga, kita dapat mengulangi pekerjaan di atas dengan
$\eusm L^1_{\Bbb C}$, ruang fungsi-fungsi terintegralkan bernilai kompleks,
sebagai pengganti $\eusm L^1$, untuk membangun ruang Banach kompleks
$L^1_{\Bbb C}$. Perubahan yang diperlukan, yang didasarkan pada
gagasan-gagasan 241J, hanya kecil.

\medskip

\noindent}{\bf (a)}\dvro{ Untuk}{ Dalam 242Aa, barangkali berguna untuk
menyatakan bahwa, untuk} $f\in\eusm L^0_{\Bbb C}$,

\Centerline{$f\in\eusm L^1_{\Bbb C}\iff |f|\in\eusm L^1\iff
\Real(f),\,\Imag(f)\in\eusm L^1$.}

\noindent Akibatnya, untuk $u\in L^0_{\Bbb C}$,

\Centerline{$u\in L^1_{\Bbb C}\iff |u|\in L^1\iff
\Real(u),\,\Imag(u)\in
L^1$.}

\header{242Pb}{\bf (b)}\cmmnt{ Untuk membuktikan versi kompleks dari
242E, perhatikan bahwa jika $\sequencen{f_n}$ suatu barisan dalam
$\eusm L^1_{\Bbb C}$ sedemikian sehingga
$\sum_{n=0}^{\infty}\int|f_n|<\infty$, maka
$\sum_{n=0}^{\infty}\int|\Real(f_n)|$ dan
$\sum_{n=0}^{\infty}\int|\Imag(f_n)|$ keduanya berhingga, sehingga kita
dapat menerapkan 242E dua kali dan melihat bahwa

\Centerline{$\int(\sum_{n=0}^{\infty}f_n)
=\int(\sum_{n=0}^{\infty}\Real(f_n))
  +i\int(\sum_{n=0}^{\infty}\Imag(f_n))
=\sum_{n=0}^{\infty}\int f_n$.}

\noindent Oleh karena itu, dengan argumen 242F kita dapat membuktikan
bahwa} $L^1_{\Bbb C}$ lengkap terhadap $\|\,\|_1$\cmmnt{ }.

\cmmnt{\header{242Pc}{\bf (c)} Demikian pula, hanya sedikit perubahan
diperlukan untuk menyesuaikan 242J agar memberikan deskripsi operator
ekspektasi bersyarat
$P:L^1_{\Bbb C}(\mu)\to L^1_{\Bbb C}(\mu\restrp\Tau)$ ketika
$(X,\Sigma,\mu)$ suatu ruang probabilitas dan $\Tau$ suatu
subaljabar-$\sigma$ dari $\Sigma$. Dalam rumus

\Centerline{$|Pu|\le P|u|$}

\noindent dari 242Je, kita perlu mengetahui bahwa

\Centerline{$|Pu|=\sup_{|\zeta|=1}\Real(\zeta Pu)$}

\noindent dalam $L^0(\mu\restrp\Tau)$ (241Jc), sedangkan

\Centerline{$\Real(\zeta Pu)=\Real(P(\zeta u))
=P(\Real(\zeta u))\le P|u|$}

\noindent kapan pun $|\zeta|=1$.
}

\cmmnt{
\header{242Pd}{\bf (d)} Dalam 242M, kita perlu mengganti $\eusm S$
dengan $\eusm S_{\Bbb C}$, yaitu ruang `fungsi sederhana bernilai kompleks'
berbentuk $\sum_{k=0}^na_k\chi E_k$, dengan setiap $a_k$ bilangan kompleks
dan setiap $E_k$ himpunan terukur berukuran hingga; kemudian kita memperoleh
subruang linear padat
$S_{\Bbb C}=\{f^{\ssbullet}:f\in\eusm S_{\Bbb C}\}$ dari
$L^1_{\Bbb C}$. Dalam 242O, kita harus mengganti $C_k$ dengan
$C_k(\BbbR^r;\Bbb C)$, yaitu ruang fungsi kontinu terbatas bernilai
kompleks yang mempunyai dukungan terbatas, dan $S_0$ dengan rentang linear
atas $\Bbb C$ dari
$\{\chi I:I$ suatu interval setengah terbuka terbatas$\}$.
}

\exercises{
\leader{242X}{Latihan dasar $\pmb{>}$(a)} Misalkan $X$ suatu himpunan,
dan misalkan $\mu$ ukuran cacah pada $X$. Tunjukkan bahwa $L^1(\mu)$
dapat diidentifikasi dengan ruang $\ell^1(X)$ dari fungsi-fungsi bernilai
riil pada $X$ yang dapat dijumlahkan secara mutlak (lihat 226A). Khususnya,
ruang $\ell^1=\ell^1(\Bbb N)$ dari barisan bernilai riil yang dapat
dijumlahkan secara mutlak merupakan suatu ruang $L^1$. Tuliskan bukti 242F
yang disesuaikan dengan kasus-kasus khusus ini.
%242F, but put at start

\sqheader 242Xb Misalkan $(X,\Sigma,\mu)$ sebarang ruang ukur, dan
$\hat\mu$ kelengkapan dari $\mu$. Tunjukkan bahwa
$\eusm L^1(\hat\mu)=\eusm L^1(\mu)$ dan
$L^1(\hat\mu)=L^1(\mu)$ (bandingkan 241Xb).
%242A

\spheader 242Xc Misalkan
$\langle(X_i,\Sigma_i,\mu_i)\rangle_{i\in I}$ suatu keluarga ruang ukur,
dan $(X,\Sigma,\mu)$ jumlah langsungnya. Tunjukkan bahwa isomorfisme antara
$L^0(\mu)$ dan $\prod_{i\in I}L^0(\mu_i)$ (241Xd) mengimbaskan suatu
identifikasi antara $L^1(\mu)$ dan

\Centerline{$\{u:u\in\prod_{i\in I}L^1(\mu_i),\,
  \|u\|=\sum_{i\in I}\|u(i)\|_1<\infty\}
\subseteq\prod_{i\in I}L^1(\mu_i)$.}
%242D

\spheader 242Xd Misalkan $(X,\Sigma,\mu)$ dan $(Y,\Tau,\nu)$
ruang-ruang ukur, dan $\phi:X\to Y$ suatu \imp\ fungsi. Tunjukkan bahwa
$g\mapsto g\phi:\eusm L^1(\nu)\to\eusm L^1(\mu)$ (235G) mengimbaskan
operator linear $T:L^1(\nu)\to L^1(\mu)$ sedemikian sehingga
$\|Tv\|_1=\|v\|_1$ untuk setiap $v\in L^1(\nu)$.
%242D

\spheader 242Xe Misalkan $U$ suatu ruang Riesz (definisi: 241Ed).
Suatu {\bf norma Riesz} pada $U$ adalah norma $\|\,\|$ sedemikian
sehingga $\|u\|\le\|v\|$ kapan pun $|u|\le|v|$. Tunjukkan bahwa jika
$U$ diberi topologi normanya (2A4Bb) untuk norma semacam itu, maka
(i) $u\mapsto|u|:U\to U$, $(u,v)\mapsto u\vee v:U\times U\to U$
kontinu (ii) $\{u:u\ge 0\}$ tertutup.
%242D

\spheader 242Xf Tunjukkan bahwa setiap kisi Banach harus merupakan
ruang Riesz Archimedes (241Fa).
%242G

\spheader 242Xg Misalkan $(X,\Sigma,\mu)$ suatu ruang probabilitas,
$\Tau$ suatu subaljabar-$\sigma$ dari $\Sigma$, dan $\Upsilon$ suatu
subaljabar-$\sigma$ dari $\Tau$. Misalkan
$P_1:L^1(\mu)\to L^1(\mu\restrp\Tau)$,
$P_2:L^1(\mu\restrp\Tau)\to L^1(\mu\restrp\Upsilon)$, dan
$P:L^1(\mu)\to L^1(\mu\restrp\Upsilon)$ operator-operator ekspektasi
bersyarat yang bersesuaian. Tunjukkan bahwa $P=P_2P_1$.
%242J

\spheader 242Xh Tunjukkan bahwa jika $g:\BbbR^r\to\Bbb R$ kontinu dan
mempunyai dukungan terbatas, maka fungsi itu terbatas dan mencapai batas-batasnya.
\Hint{2A2F-2A2G.}
%242O

\spheader 242Xi Misalkan $\mu$ ukuran Lebesgue pada $\Bbb R$. (i) Ambil
$\delta>0$. Tunjukkan bahwa jika
$\phi_{\delta}(x)=\exp(-\Bover{1}{\delta^2-x^2})$ untuk
$|x|<\delta$, dan $0$ untuk $|x|\ge\delta$, maka $\phi_{\delta}$ {\bf mulus},
yakni dapat didiferensialkan sebanyak kali yang diinginkan. (ii) Tunjukkan
bahwa jika $F_{\delta}(x)=\int_{-\infty}^x\phi_{\delta}d\mu$ untuk
$x\in\Bbb R$, maka $F_{\delta}$ mulus. (iii) Tunjukkan bahwa jika
$a<b<c<d$ dalam $\Bbb R$, terdapat fungsi mulus $h$ sedemikian sehingga
$\chi[b,c]\le h\le\chi[a,d]$. (iv) Tuliskan $\eusm D$ untuk ruang
fungsi-fungsi mulus $h:\Bbb R\to\Bbb R$ sedemikian sehingga
$\{x:h(x)\ne 0\}$ terbatas. Tunjukkan bahwa
$\{h^{\ssbullet}:h\in\eusm D\}$ padat dalam $L^1(\mu)$.
(v) Misalkan $f$ suatu fungsi bernilai riil yang terintegralkan pada setiap
subhimpunan terbatas dari $\Bbb R$. Tunjukkan bahwa $f\times h$
terintegralkan untuk setiap $h\in\eusm D$, dan bahwa jika
$\int f\times h=0$ untuk setiap $h\in\eusm D$, maka $f=0$ hampir di
mana-mana. \Hint{222D.}
%242O

\spheader 242Xj Misalkan $(X,\Sigma,\mu)$ suatu ruang probabilitas,
$\Tau$ suatu subaljabar-$\sigma$ dari $\Sigma$, dan
$P:L^1(\mu)\to L^1(\mu\restrp\Tau)\subseteq L^1(\mu)$ operator
ekspektasi bersyarat yang bersesuaian. Tunjukkan bahwa jika $u$,
$v\in L^1(\mu)$ sedemikian sehingga $P|u|\times P|v|\in L^1(\mu)$,
maka $\int Pu\times v=\int Pu\times Pv=\int u\times Pv$.
%242L

\leader{242Y}{Latihan lanjutan (a)}
%\spheader 242Ya
Misalkan $(X,\Sigma,\mu)$ suatu ruang ukur. Misalkan
$A\subseteq L^1=L^1(\mu)$ suatu himpunan tak kosong yang terarah ke bawah,
dan andaikan bahwa $\inf A=0$ dalam $L^1$. (i) Tunjukkan bahwa
$\inf_{u\in A}\|u\|_1=0$.
\Hint{tetapkan $\gamma=\inf_{u\in A}\|u\|_1$; temukan suatu barisan tak
menaik $\sequencen{u_n}$ dalam $A$ sedemikian sehingga
$\lim_{n\to\infty}\|u_n\|_1=\gamma$; tetapkan
$v=\inf_{n\in\Bbb N}u_n$ dan tunjukkan bahwa $u\wedge v=v$ untuk setiap
$u\in A$, sehingga $v=0$.} (ii) Tunjukkan bahwa jika $U$ sebarang
himpunan terbuka yang memuat $0$, terdapat $u\in A$ sedemikian sehingga
$v\in U$ kapan pun $0\le v\le u$.
%242D

\spheader 242Yb Misalkan $(X,\Sigma,\mu)$ suatu ruang ukur dan $Y$
sebarang subhimpunan dari $X$; misalkan $\mu_Y$ ukuran subruang pada $Y$
dan $T:L^0(\mu)\to L^0(\mu_Y)$ peta kanonik yang dideskripsikan dalam
241Yg. (i) Tunjukkan bahwa $Tu\in L^1(\mu_Y)$ dan
$\|Tu\|_1\le\|u\|_1$ untuk setiap $u\in L^1(\mu)$. (ii) Tunjukkan
bahwa jika $u\in L^1(\mu)$, maka $\|Tu\|_1=\|u\|_1$ jika dan hanya jika
$\int_Eu=\int_{Y\cap E}Tu$ untuk setiap $E\in\Sigma$. (iii) Tunjukkan
bahwa $T$ surjektif dan bahwa
$\|v\|_1=\min\{\|u\|_1:u\in L^1(\mu),\,Tu=v\}$ untuk setiap
$v\in L^1(\mu_Y)$. \Hint{214Eb.} (Lihat juga 244Yd di bawah.)
%242D

\spheader 242Yc Misalkan $(X,\Sigma,\mu)$ suatu ruang ukur. Tuliskan
$\eusm L^1_{\Sigma}$ untuk ruang semua fungsi $\Sigma$-terukur yang
terintegralkan dari $X$ ke $\Bbb R$, dan $\eusm N$ untuk subruang dari
$\eusm L^1_{\Sigma}$ yang terdiri atas fungsi-fungsi terukur yang bernilai
nol hampir di mana-mana. (i) Tunjukkan bahwa $\eusm L^1_{\Sigma}$
merupakan ruang Riesz lengkap-$\sigma$ Dedekind. (ii) Tunjukkan bahwa
$L^1(\mu)$ dapat diidentifikasi, sebagai ruang linear terurut, dengan hasil
bagi $\eusm L^1_{\Sigma}/\eusm N$ sebagaimana didefinisikan dalam 241Yb.
(iii) Tunjukkan bahwa $\|\,\|_1$ merupakan seminorma pada
$\eusm L^1_{\Sigma}$. (iv) Tunjukkan bahwa
$f\mapsto|f|:\eusm L^1_{\Sigma}\to\eusm L^1_{\Sigma}$ kontinu jika
$\eusm L^1_{\Sigma}$ diberi topologi yang didefinisikan dari
$\|\,\|_1$. (v) Tunjukkan bahwa $\{f:f=0$ hampir di mana-mana$\}$
tertutup dalam $\eusm L^1_{\Sigma}$, tetapi $\{f:f\ge 0\}$ belum tentu
tertutup.
%242D

\spheader 242Yd Misalkan $(X,\Sigma,\mu)$ suatu ruang ukur, dan
$\tilde\mu$ versi c.l.d.\ dari $\mu$ (213E). Tunjukkan bahwa inklusi
$\eusm L^1(\mu)\subseteq\eusm L^1(\tilde\mu)$ mengimbaskan suatu
isomorfisme, sebagai ruang linear bernorma terurut, antara
$L^1(\tilde\mu)$ dan $L^1(\mu)$.
%242D

\spheader 242Ye Misalkan $(X,\Sigma,\mu)$ suatu ruang ukur dan
$u_0,\ldots,u_n\in L^1(\mu)$. (i) Andaikan $k_0,\ldots,k_n\in\Bbb Z$
sedemikian sehingga $\sum_{i=0}^nk_i=1$. Tunjukkan bahwa
$\sum_{i=0}^n\sum_{j=0}^nk_ik_j\|u_i-u_j\|_1\le 0$. \Hint{
$\sum_{i=0}^n\sum_{j=0}^n
\ifdim\pagewidth=390pt\penalty-100\fi
k_ik_j|\alpha_i-\alpha_j|\le 0$ untuk semua
$\alpha_0,\ldots,\alpha_n\in\Bbb R$.} (ii) Andaikan
$\gamma_0,\ldots,\gamma_n\in\Bbb R$ sedemikian sehingga
$\sum_{i=0}^n\gamma_i=0$. Tunjukkan bahwa
$\sum_{i=0}^n\sum_{j=0}^n\gamma_i\gamma_j\|u_i-u_j\|_1\le 0$.
%242D

\spheader 242Yf Misalkan $(X,\Sigma,\mu)$ suatu ruang ukur, dan
$A\subseteq L^1=L^1(\mu)$ suatu himpunan tak kosong yang terarah ke atas.
Andaikan bahwa $A$ terbatas terhadap norma $\|\,\|_1$. (i) Tunjukkan
bahwa terdapat barisan tak menurun $\sequencen{u_n}$ dalam $A$ sedemikian
sehingga $\lim_{n\to\infty}\int u_n=\sup_{u\in A}\int u$, dan bahwa
$\sequencen{u_n}$ merupakan barisan Cauchy. (ii) Tunjukkan bahwa
$w=\sup A$ terdefinisi dalam $L^1$ dan termasuk dalam penutupan menurut
norma dari $A$ dalam $L^1$, sehingga, khususnya,
$\|w\|_1\le\sup_{u\in A}\|u\|_1$.
%242F

\spheader 242Yg Suatu norma Riesz (definisi: 242Xe) pada ruang Riesz
$U$ disebut {\bf kontinu menurut urutan} jika
$\inf_{u\in A}\|u\|=0$ kapan pun $A\subseteq U$ suatu himpunan tak
kosong yang terarah ke bawah dengan infimum $0$. (Jadi 242Ya menyatakan
bahwa semua norma $\|\,\|_1$ kontinu menurut urutan.) Tunjukkan bahwa
dalam hal ini (i) setiap barisan tak menurun dalam $U$ yang mempunyai batas
atas dalam $U$ harus merupakan barisan Cauchy (ii) jika $U$ suatu kisi
Banach, maka $U$ lengkap Dedekind. ({\it Petunjuk untuk (i)\/}: jika
$\sequencen{u_n}$ suatu barisan tak menurun dengan batas atas dalam $U$,
misalkan $B$ himpunan batas atas dari $\{u_n:n\in\Bbb N\}$ dan tunjukkan
bahwa $A=\{v-u_n:v\in B,\,n\in\Bbb N\}$ mempunyai infimum $0$ karena
$U$ bersifat Archimedes.)
%242Yf 242G

\spheader 242Yh Misalkan $(X,\Sigma,\mu)$ sebarang ruang ukur.
Tunjukkan bahwa $L^1(\mu)$ mempunyai sifat supremum terhitung (241Ye).
%242Yg 242Yf 242G

\spheader 242Yi Secara lebih umum, tunjukkan bahwa setiap ruang Riesz
dengan norma Riesz yang kontinu menurut urutan mempunyai sifat supremum
terhitung.
%242Yh 242Yg 242Yf 242G

\spheader 242Yj Misalkan $(X,\Sigma,\mu)$ dan $(Y,\Tau,\nu)$
ruang-ruang ukur dan $U\subseteq L^0(\mu)$ suatu subruang linear.
Misalkan $T:U\to L^0(\nu)$ suatu operator linear sedemikian sehingga
$Tu\ge 0$ dalam $L^0(\nu)$ kapan pun $u\in U$ dan $u\ge 0$ dalam
$L^0(\mu)$. Andaikan $w\in U$ sedemikian sehingga $w\ge 0$ dan
$Tw=(\chi Y)^{\ssbullet}$. Tunjukkan bahwa kapan pun
$\phi:\Bbb R\to\Bbb R$ suatu fungsi cembung dan $u\in L^0(\mu)$
sedemikian sehingga $w\times u$ dan $w\times\bar\phi(u)\in U$, dengan
mendefinisikan $\bar\phi:L^0(\mu)\to L^0(\mu)$ seperti dalam 241I,
maka $\bar\phi T(w\times u)\le T(w\times\bar\phi u)$. Jelaskan bagaimana
hasil ini dapat dipandang sebagai perumuman bersama ketaksamaan Jensen,
sebagaimana dinyatakan dalam 233I, dan 242K di atas. Lihat juga 244M di
bawah.
%242K

\spheader 242Yk(i) Suatu fungsi $\phi:\Bbb C\to\Bbb R$ disebut
{\bf cembung} jika
$\phi(ab+(1-a)c)\le a\phi(b)+(1-a)\phi(c)$ untuk semua $b$,
$c\in\Bbb C$ dan $a\in [0,1]$. (ii) Tunjukkan bahwa fungsi semacam itu
harus terbatas pada setiap subhimpunan terbatas dari $\Bbb C$. (iii) Jika
$\phi:\Bbb C\to\Bbb R$ cembung dan $c\in\Bbb C$, tunjukkan bahwa
terdapat $b\in\Bbb C$ sedemikian sehingga
$\phi(x)\ge\phi(c)+\Real(b(x-c))$ untuk setiap $x\in\Bbb C$. (iv) Jika
$\langle b_c\rangle_{c\in\Bbb C}$ sedemikian sehingga
$\phi(x)\ge\phi_c(x)=\phi(c)+\Real(b_c(x-c))$ untuk semua
$x$, $c\in\Bbb C$, tunjukkan bahwa $\{b_c:c\in I\}$ terbatas untuk
setiap $I\subseteq\Bbb C$ terbatas. (v) Tunjukkan bahwa jika
$D\subseteq\Bbb C$ sebarang himpunan padat, maka
$\phi(x)=\sup_{c\in D}\phi_c(x)$ untuk setiap $x\in\Bbb C$.
%242P

\spheader 242Yl Misalkan $(X,\Sigma,\mu)$ suatu ruang probabilitas dan
$\Tau$ suatu subaljabar-$\sigma$ dari $\Sigma$. Misalkan
$P:L^1_{\Bbb C}(\mu)\to L^1_{\Bbb C}(\mu\restrp\Tau)$ operator
ekspektasi bersyarat. Tunjukkan bahwa jika
$\phi:\Bbb C\to\Bbb R$ sebarang fungsi cembung, dan kita mendefinisikan
$\bar\phi(f^{\ssbullet})=(\phi f)^{\ssbullet}$ untuk setiap
$f\in\eusm L^0_{\Bbb C}(\mu)$, maka
$\bar\phi(Pu)\le P(\bar\phi(u))$ kapan pun
$u\in L^1_{\Bbb C}(\mu)$ sedemikian sehingga
$\bar\phi(u)\in L^1(\mu)$.
%242P
}%end of exercises

\endnotes{
\Notesheader{242} Tentu saja ruang-ruang $L^1$ membentuk salah satu kelas
ruang Riesz terpenting, dan oleh karena itu sifat-sifatnya menempati kedudukan
utama dalam teori umum; 242Xe, 242Xf, 242Ya, dan 242Yf-242Yi %242Yf 242Yg 242Yh 242Yi
menguraikan beberapa keterkaitan di antara sifat-sifat ini. Saya akan kembali
ke pertanyaan-pertanyaan tersebut dalam Bab 35 di jilid berikutnya. Secara
sepintas saya telah menyebutkan (242Dd) aditivitas norma $L^1$ pada unsur-unsur
positif. Fakta elementer ini sesungguhnya mencirikan ruang-ruang $L^1$ di
antara kisi-kisi Banach; lihat 369E di jilid berikutnya.

Sebagaimana $L^0(\mu)$ dapat dipandang sebagai hasil bagi suatu ruang linear
$\eusm L^0_{\Sigma}$, demikian pula $L^1(\mu)$ dapat dipandang sebagai
hasil bagi suatu ruang linear $\eusm L^1_{\Sigma}$ (242Yc). Saya telah
membahas pertanyaan ini dalam catatan untuk \S241; yang saya upayakan di sini
hanyalah bersikap konsisten.

Kini kita mempunyai bahasa yang memungkinkan kita membicarakan ekspektasi
bersyarat `dari' suatu fungsi $f$, yaitu kelas ekuivalensi dalam
$L^1(\mu\restrp\Tau)$ yang terdiri tepat atas semua ekspektasi bersyarat
dari $f$ pada $\Tau$. Jika kita memandang $L^1(\mu\restrp\Tau)$ sebagai
ruang yang diidentifikasi dengan citranya dalam $L^1(\mu)$, maka operator
ekspektasi bersyarat $P:L^1(\mu)\to L^1(\mu\restrp\Tau)$ menjadi suatu
proyeksi (242Jh). Dengan demikian, kita mempunyai pernyataan ulang
233J-233K, seperti dalam 242K, 242L, dan 242Yj.

Saya memberikan 242O dalam bentuk yang cukup umum; tetapi arti pentingnya
sudah tampak jika kita mengambil $X$ sebagai $[0,1]$ dengan ukuran Lebesgue satu
dimensi. Dalam kasus ini, kita mempunyai norma alami pada $C([0,1])$, yaitu
ruang semua fungsi kontinu bernilai riil pada $[0,1]$, yang diberikan dengan
menetapkan

\Centerline{$\|f\|_1=\int_0^1|f(x)|dx$}

\noindent untuk setiap $f\in C([0,1])$. Integral di sini tentu saja dapat
diambil sebagai integral Riemann; kita tidak memerlukan teori Lebesgue untuk
menunjukkan bahwa $\|\,\|_1$ merupakan norma pada $C([0,1])$. Mudah
diperiksa bahwa $C([0,1])$ tidak lengkap terhadap norma ini (jika kita
menetapkan $f_n(x)=\min(1,2^nx^n)$ untuk $x\in[0,1]$, maka
$\sequencen{f_n}$ merupakan barisan Cauchy terhadap $\|\,\|_1$ yang tidak
mempunyai limit terhadap $\|\,\|_1$ dalam $C([0,1])$). Kita dapat
menggunakan teori abstrak ruang bernorma untuk membangun kelengkapan
$C([0,1])$; tetapi jauh lebih memuaskan jika kelengkapan ini dapat diberi
bentuk yang relatif konkret, dan inilah yang dapat dilakukan oleh
identifikasi $L^1$ dengan kelengkapan $C([0,1])$. (Perhatikan bahwa
pernyataan bahwa $\|\,\|_1$ merupakan norma pada $C([0,1])$, yakni bahwa
$\|f\|_1\ne 0$ untuk setiap $f\in C([0,1])$ yang tak nol, berarti tepat
bahwa peta $f\mapsto f^{\ssbullet}:C([0,1])\to L^1$ bersifat injektif,
sehingga $C([0,1])$ dapat diidentifikasi, sebagai ruang bernorma terurut,
dengan citranya dalam $L^1$.) Akan lebih baik lagi jika kita dapat menemukan
suatu realisasi kelengkapan $C([0,1])$ sebagai ruang fungsi pada suatu
himpunan $Z$, alih-alih sebagai ruang kelas ekuivalensi fungsi pada $[0,1]$.
Sayangnya hal ini tidak praktis; realisasi semacam itu memang ada, tetapi
harus melibatkan entah suatu himpunan dasar $Z$ yang sama sekali tidak
familier, atau suatu peta pembenaman dari $C([0,1])$ ke $\BbbR^Z$ yang
terlampau sewenang-wenang.

Anda dapat memperoleh gambaran tentang rintangan dalam merealisasikan
kelengkapan $C([0,1])$ sebagai ruang fungsi pada $[0,1]$ itu sendiri dengan
meninjau $f_n(x)=\bover1nx^n$ untuk $n\ge 1$. Perhitungan mudah menunjukkan
bahwa $\sum_{n=1}^{\infty}\|f_n\|_1<\infty$, sehingga
$\sum_{n=1}^{\infty}f_n$ harus ada dalam kelengkapan $C([0,1])$; tetapi tidak
ada nilai alami yang dapat diberikan kepadanya pada titik $1$. Penyesuaian
gagasan ini dapat menimbulkan gejala yang rumit tanpa batas -- bahkan, 242O
menunjukkan bahwa setiap fungsi terintegralkan berhubungan dengan suatu
barisan yang sesuai dari $C([0,1])$. Dalam \S245 saya akan membahas lebih
lanjut seperti apa barisan yang konvergen terhadap $\|\,\|_1$.

Dari sudut pandang teori ukuran yang dipahami secara sempit, sebagian besar
gagasan menarik tampak paling jelas dengan fungsi riil dan ruang linear riil.
Namun, beberapa penerapan terpenting teori ukuran -- penting bukan hanya bagi
matematika secara umum, tetapi juga bagi pertanyaan teori-ukur yang
diilhaminya -- berkenaan dengan fungsi kompleks dan ruang linear kompleks.
Oleh karena itu, saya terus menawarkan ikhtisar teori kompleks, seperti dalam
242P. Saya mencatat bahwa pada selang yang tidak teratur kita memerlukan
gagasan-gagasan yang belum dijabarkan dalam teori riil, seperti dalam 242Pb
dan 242Yl.
}%end of notes


\discrpage
